\chapter*{Đề cương}
\label{proposal}

\section*{1. Thông tin chung}

\noindent
\textbf{Tên đề tài:} Nghiên cứu phương pháp trả lời câu hỏi có ràng buộc về thời gian sử dụng suy
luận trừu tượng (\textit{Abstract Reasoning for Temporal Question Answering}).

\medskip
\noindent
\textbf{Người hướng dẫn:}
\begin{itemize}
  \item TS. Nguyễn Hồng Bửu Long (Khoa Công nghệ Thông tin).
  \item TS. Lương An Vinh (Trường Đại học Công nghệ Sài Gòn).
\end{itemize}

\noindent
\textbf{Sinh viên thực hiện:} Liêu Hải Lưu Danh (MSSV: 22120459).\\
\textbf{Loại đề tài:} Nghiên cứu.\\
\textbf{Thời gian thực hiện:} Từ 02/2026 đến 07/2026.

\section*{2. Nội dung thực hiện}

\subsection*{2.1. Giới thiệu về đề tài}

Sự phát triển của các mô hình ngôn ngữ lớn (LLMs) đã thúc đẩy mạnh mẽ các hệ thống Hỏi--Đáp tri
thức. Tuy nhiên, các mô hình này vẫn gặp khó khăn trong những bài toán yêu cầu suy luận nhiều bước,
đặc biệt trên đồ thị tri thức có yếu tố thời gian, nơi thông tin không chỉ cần được truy xuất mà còn
phải được xử lý theo trình tự logic.

Đề tài tập trung nghiên cứu và áp dụng phương pháp Suy luận Trừu tượng Quy nạp
(\textit{Abstract Reasoning Induction} -- ARI)~\cite{chen2024ari} nhằm nâng cao năng lực suy luận
cho hệ thống Hỏi--Đáp tri thức thời gian. Ý tưởng chính là học các chiến lược suy luận trừu tượng
từ những bài toán đã giải trước đó, sau đó sử dụng các chiến lược này để định hướng quá trình giải
quyết câu hỏi mới. Hệ thống được xây dựng theo hai thành phần: khối xử lý dựa trên tri thức để truy
vấn và thao tác trên đồ thị tri thức, và khối suy luận sử dụng mô hình ngôn ngữ lớn để lựa chọn
hành động phù hợp dựa trên hướng dẫn suy luận đã học.

Ví dụ, với câu hỏi \textit{``Tổng thống Hoa Kỳ khi Chiến tranh Việt Nam kết thúc là ai?''}, hệ
thống cần thực hiện nhiều bước suy luận: xác định thời điểm kết thúc Chiến tranh Việt Nam (1975),
sau đó truy vấn ai là tổng thống Hoa Kỳ tại thời điểm đó. Những câu hỏi dạng này yêu cầu kết hợp
thông tin thực thể, quan hệ và mốc thời gian trong đồ thị tri thức.

\subsection*{2.2. Mục tiêu đề tài}

Đề tài được thực hiện nhằm giải quyết các hạn chế trên thông qua việc áp dụng phương pháp ARI, cho
phép hệ thống học và tái sử dụng các chiến lược suy luận trừu tượng từ dữ liệu lịch sử. Cách tiếp
cận này giúp chuyển từ suy luận tức thời sang suy luận có định hướng, từ đó nâng cao hiệu quả giải
quyết các câu hỏi phức tạp. Kết quả của đề tài góp phần cải thiện năng lực suy luận nhiều bước cho
các hệ thống Hỏi--Đáp tri thức thời gian, phục vụ các ứng dụng thực tiễn như tra cứu sự kiện lịch
sử, tiểu sử nhân vật và phân tích dữ liệu theo dòng thời gian.

\subsection*{2.3. Phạm vi của đề tài}

\textbf{Về lý thuyết:} đề tài nghiên cứu và ứng dụng phương pháp suy luận trừu tượng ARI cho bài
toán hỏi--đáp trên đồ thị tri thức có yếu tố thời gian (\textit{Temporal Knowledge Graph Question
Answering}), kết hợp khả năng suy luận của mô hình ngôn ngữ lớn với các thao tác truy xuất tri thức
có cấu trúc.

\textbf{Về đối tượng nghiên cứu:} đề tài tập trung vào các câu hỏi yêu cầu suy luận nhiều bước trên
đồ thị con của đồ thị tri thức thời gian~\cite{chen2022subgraph}, đồng thời tham khảo các hướng kết
hợp mô hình ngôn ngữ lớn với tri thức để hỏi--đáp trên đồ thị tri thức~\cite{baek2023kaping}.

\textbf{Về nguồn dữ liệu:} đề tài sử dụng các bộ dữ liệu hỏi--đáp trên đồ thị tri thức thời gian đã
công bố như MultiTQ~\cite{chen2023multitq} và CronQuestions~\cite{saxena2021cronkgqa} để tham chiếu,
huấn luyện và đánh giá.

\textbf{Về mô hình và hệ thống:} hệ thống được xây dựng theo kiến trúc hai pha gồm: (1) pha dựa trên
tri thức để truy xuất và xử lý dữ liệu từ đồ thị, và (2) pha suy luận trừu tượng dùng mô hình ngôn
ngữ lớn để lựa chọn và thực thi các hành động suy luận.

\subsection*{2.4. Cách tiếp cận dự kiến}

Hệ thống được xây dựng theo hai giai đoạn chính.

\textbf{Giai đoạn thứ nhất (Knowledge-based Phase)} tập trung vào việc xử lý dữ liệu và tương tác
với đồ thị tri thức thời gian: truy vấn, lọc và trích xuất thông tin liên quan dựa trên câu hỏi đầu
vào, đồng thời ghi nhận các bước suy luận và kết quả trung gian.

\textbf{Giai đoạn thứ hai (Abstract Reasoning Phase)} sử dụng mô hình ngôn ngữ lớn để tổng quát hóa
các chuỗi hành động đã thực hiện thành các chiến lược suy luận trừu tượng. Các chiến lược này được
lưu trữ và sử dụng làm hướng dẫn khi hệ thống xử lý các câu hỏi mới, giúp lựa chọn hành động suy
luận phù hợp.

Cụ thể, quá trình giải một câu hỏi gồm các bước: (1) tiền xử lý câu hỏi để xác định thực thể, quan
hệ và yếu tố thời gian; (2) truy xuất tri thức từ đồ thị; (3) thực hiện chuỗi hành động suy luận;
(4) trích xuất chiến lược suy luận trừu tượng từ các chuỗi hành động; (5) lưu trữ và tái sử dụng
chiến lược; (6) áp dụng chiến lược phù hợp cho câu hỏi mới để xác định câu trả lời cuối cùng.

\subsection*{2.5. Kết quả dự kiến của đề tài}

\textbf{Về kết quả định lượng:} hệ thống được kỳ vọng cải thiện độ chính xác trong bài toán hỏi--đáp
trên đồ thị tri thức có yếu tố thời gian so với các phương pháp suy luận trực tiếp bằng mô hình ngôn
ngữ lớn. Các chỉ số đánh giá dự kiến bao gồm độ chính xác của câu trả lời, khả năng suy luận nhiều
bước và hiệu quả truy xuất thông tin.

\textbf{Về sản phẩm đầu ra:} đề tài dự kiến xây dựng một hệ thống thử nghiệm cho bài toán hỏi--đáp
trên đồ thị tri thức có yếu tố thời gian, áp dụng phương pháp ARI, có khả năng nhận câu hỏi đầu vào,
truy xuất thông tin, thực hiện suy luận nhiều bước và trả về câu trả lời tương ứng.

\subsection*{2.6. Kế hoạch thực hiện}

\begin{table}[ht]
\centering
\begin{tabular}{C{2.2cm} L{10.3cm}}
\toprule
\textbf{Thời gian} & \textbf{Kế hoạch} \\
\midrule
02/2026 & Khảo sát tài liệu liên quan đến hệ thống hỏi--đáp trên đồ thị tri thức, nghiên cứu phương pháp Abstract Reasoning Induction và các mô hình suy luận nhiều bước. \\
03/2026 & Phân tích bài toán, xác định phạm vi nghiên cứu và thiết kế kiến trúc hệ thống cho bài toán hỏi--đáp trên Temporal Knowledge Graph. \\
04/2026 & Xây dựng pipeline xử lý câu hỏi, cài đặt các mô-đun truy xuất tri thức và các thao tác suy luận cơ bản trên đồ thị tri thức. \\
05/2026 & Tích hợp mô hình ngôn ngữ lớn để thực hiện suy luận trừu tượng và xây dựng cơ chế lưu trữ, tái sử dụng chiến lược suy luận. \\
06/2026 & Thực nghiệm hệ thống trên bộ dữ liệu đã chọn, đánh giá kết quả và so sánh với các phương pháp liên quan. \\
07/2026 & Phân tích kết quả, hoàn thiện hệ thống, viết báo cáo khóa luận và chuẩn bị trình bày kết quả nghiên cứu. \\
\bottomrule
\end{tabular}
\end{table}
